\documentclass[11pt]{article}

\usepackage{amsmath,amssymb,amsthm,mathtools}
\usepackage[margin=1in]{geometry}
\usepackage{microtype}

\newtheorem{theorem}{Theorem}[section]
\newtheorem{lemma}[theorem]{Lemma}
\newtheorem{corollary}[theorem]{Corollary}
\newtheorem{remark}[theorem]{Remark}

\newcommand{\dd}{\,\mathrm d}

\title{A Triangle with a Two-Leaf Broom:\\
A Rooted Supporting-Plane Inequality}
\author{}
\date{}

\begin{document}

\maketitle

\begin{abstract}
Let $B$ be obtained from a triangle by joining one triangle vertex to the
centre of a two-leaf star.  We prove that every graphon $W$ of edge density
$p\geq1/2$ satisfies
\[
 t(B,W)\geq p^4(2p-1)=\Phi_B(p).
\]
The proof is edge-density-sensitive.  It does not use the false universal
comparison between the associated five-edge tree and
$S_4\sqcup2K_2$.  Instead, two pointwise graphon inequalities reduce the
problem to a sharp two-variable scalar inequality on
\[
 0\leq a\leq d,\qquad a\geq d-(1-p).
\]
The scalar lower bound contains a multiple of $a-d^2$, whose integral is
zero for $a=T_Wd$.  This proves the formerly open NetworkX Atlas graph
$100$.
\end{abstract}


%%%%%%%%%%%%%%%%%%%%%%%%%%%%%%%%%%%%%%%%%%%%%%%%%%%%%%%%%%%%%%%%%%%%%%%%%%%%%%%
\section{The graph, its target, and rooted quantities}
%%%%%%%%%%%%%%%%%%%%%%%%%%%%%%%%%%%%%%%%%%%%%%%%%%%%%%%%%%%%%%%%%%%%%%%%%%%%%%%

Let $(\Omega,\mu)$ be an arbitrary probability space and let
$W\colon\Omega^2\to[0,1]$ be symmetric and measurable.  Put
\[
 p=\int_{\Omega^2}W(x,y)\dd\mu(x)\dd\mu(y),
 \qquad
 d(x)=\int_\Omega W(x,y)\dd\mu(y),
\]
and define
\[
 A=T_Wd,
 \qquad
 C=T_W(d^2).
\]
The rooted triangle density is
\[
 \tau(x)=\int_{\Omega^2}
 W(x,y)W(x,z)W(y,z)\dd\mu(y)\dd\mu(z).
\]

Construct $B$ from a triangle rooted at $o$ by adding three new vertices
$u,v,w$ and exactly the three further edges
\[
 ou,\qquad uv,\qquad uw.
\]
After fixing the image $x$ of $o$, integrating $v,w$, and then $u$ gives
the exact identity
\begin{equation}
 t(B,W)=\int_\Omega \tau(x)C(x)\dd\mu(x).
 \label{eq:density}
\end{equation}

The graph has six vertices and six edges.  After properly colouring the
triangle, each successive vertex of the attached tree has $x-1$ choices.
Therefore
\[
 \chi_B(x)=x(x-1)^4(x-2),
 \qquad \chi(B)=3.
\]
For $q=1-p>0$,
\[
 q^6\chi_B(1/q)=p^4(2p-1).
\]
Thus the polynomial continuation of the target is
\begin{equation}
 \Phi_B(p)=p^4(2p-1),
 \label{eq:target}
\end{equation}
and the complete required interval is $p\in[1/2,1]$.


%%%%%%%%%%%%%%%%%%%%%%%%%%%%%%%%%%%%%%%%%%%%%%%%%%%%%%%%%%%%%%%%%%%%%%%%%%%%%%%
\section{Two pointwise graphon reductions}
%%%%%%%%%%%%%%%%%%%%%%%%%%%%%%%%%%%%%%%%%%%%%%%%%%%%%%%%%%%%%%%%%%%%%%%%%%%%%%%

\begin{lemma}[Rooted feasible region]\label{lem:region}
For almost every $x$, writing $d=d(x)$ and $a=A(x)$, one has
\[
 0\leq a\leq d,
 \qquad
 a\geq d-(1-p).
\]
Moreover,
\[
 \int_\Omega A\dd\mu=\int_\Omega d^2\dd\mu.
\]
\end{lemma}

\begin{proof}
Nonnegativity is immediate.  Since $d(y)\leq1$ almost everywhere,
\[
 A(x)=\int W(x,y)d(y)\dd\mu(y)\leq d(x).
\]
The elementary inequality $rs\geq r+s-1$ for $r,s\in[0,1]$, applied to
$r=W(x,y)$ and $s=d(y)$, gives
\[
 A(x)\geq\int\bigl(W(x,y)+d(y)-1\bigr)\dd\mu(y)
 =d(x)+p-1.
\]
Finally Tonelli's theorem and symmetry give
\[
 \int A(x)\dd\mu(x)
 =\int_{\Omega^2}W(x,y)d(y)\dd\mu^2
 =\int_\Omega d(y)^2\dd\mu(y).
\]
\end{proof}

\begin{lemma}[Broom and triangle roots]\label{lem:pointwise}
For almost every $x$,
\[
 C(x)\geq
 \begin{cases}
 A(x)^2/d(x),&d(x)>0,\\
 0,&d(x)=0,
 \end{cases}
 \qquad
 \tau(x)\geq\bigl(2A(x)-p\bigr)_+.
\]
Here $r_+=\max\{r,0\}$.
\end{lemma}

\begin{proof}
For $d(x)>0$, Cauchy--Schwarz with respect to the finite measure
$W(x,y)\dd\mu(y)$ gives
\[
 A(x)^2
 =\left(\int W(x,y)d(y)\dd\mu(y)\right)^2
 \leq d(x)\int W(x,y)d(y)^2\dd\mu(y)
 =d(x)C(x).
\]
If $d(x)=0$, then $W(x,y)=0$ for almost every $y$, and hence
$A(x)=C(x)=0$.

For the triangle estimate, apply $rs\geq r+s-1$ to
$W(x,y),W(x,z)$, multiply by $W(y,z)\geq0$, and integrate in $y,z$.
This gives $\tau(x)\geq2A(x)-p$.  Also $\tau(x)\geq0$.
\end{proof}


%%%%%%%%%%%%%%%%%%%%%%%%%%%%%%%%%%%%%%%%%%%%%%%%%%%%%%%%%%%%%%%%%%%%%%%%%%%%%%%
\section{The sharp scalar supporting inequality}
%%%%%%%%%%%%%%%%%%%%%%%%%%%%%%%%%%%%%%%%%%%%%%%%%%%%%%%%%%%%%%%%%%%%%%%%%%%%%%%

The next lemma is the new analytic input.  Its term $a-d^2$ is tailored to
the last identity in Lemma~\ref{lem:region}.

\begin{lemma}[Dense rooted supporting plane]\label{lem:scalar}
Let $p\in[1/2,1]$, and let $d,a\in[0,1]$ satisfy
\[
 0\leq a\leq d,
 \qquad a\geq d-(1-p).
\]
Define the left side to be zero when $d=0$.  Then
\begin{align}
 \frac{a^2}{d}(2a-p)_+
 \geq{}&2p^4(1-4p)+p^3(10p-3)d
 \notag\\
 &\quad+2p^2(3p-1)(a-d^2).
 \label{eq:scalar}
\end{align}
\end{lemma}

\begin{proof}
Write the right side of \eqref{eq:scalar} as
\[
 L_p(d,a)=\alpha+\beta d+\gamma(a-d^2),
\]
where
\[
 \alpha=2p^4(1-4p),\qquad
 \beta=p^3(10p-3),\qquad
 \gamma=2p^2(3p-1)>0.
\]
If $d=0$, then $a=0$, while $L_p(0,0)=\alpha\leq0$.
We may therefore assume $d>0$.

First suppose $a\leq p/2$.  The left side of \eqref{eq:scalar} is zero,
and $L_p(d,a)$ is increasing in $a$.  If $d\leq p/2$, it is enough to put
$a=d$.  Direct expansion gives
\[
 L_p(d,d)=-p^2Q_1(d),
\]
where
\[
 Q_1(d)=2(3p-1)d^2+(-10p^2-3p+2)d+8p^3-2p^2.
\]
The derivative of $Q_1$ at $d=p/2$ is
$-4p^2-5p+2<0$.  Since $Q_1$ is convex, it is decreasing on
$[0,p/2]$.  Hence
\[
 Q_1(d)\geq Q_1(p/2)
 =\frac p2(9p^2-8p+2)>0;
\]
the last quadratic has negative discriminant.

If instead $d\geq p/2$, it is enough to put $a=p/2$.  Feasibility also
forces $d\leq1-p/2$.  Here
\[
 L_p(d,p/2)=-p^2Q_2(d),
\]
where
\[
 Q_2(d)=2(3p-1)d^2+(-10p^2+3p)d+8p^3-5p^2+p.
\]
Its leading coefficient is positive and its discriminant is
\[
 -p(2p-1)^2(23p-8)\leq0.
\]
Thus $Q_2(d)\geq0$.  This proves \eqref{eq:scalar} when $a\leq p/2$.

Now suppose $a\geq p/2$.  Multiplication by $d>0$ reduces the assertion to
\begin{equation}
 G_p(d,a):=2a^3-pa^2-dL_p(d,a)\geq0.
 \label{eq:G}
\end{equation}
For fixed $d$,
\[
 \frac{\partial^2G_p}{\partial a^2}=12a-2p\geq4p>0.
\]
Consequently the minimum on the feasible interval
\[
 \max\{p/2,d-(1-p)\}\leq a\leq d
\]
occurs at an endpoint or at the unique stationary point.  We check all
possibilities exactly.

At $a=p/2$, expansion gives
\[
 G_p(d,p/2)=dp^2Q_2(d)\geq0.
\]
At $a=d$, write $G_p(d,d)=dP_p(d)$, where
\[
 P_p(d)=2(3p^3-p^2+1)d^2
 +(-10p^4-3p^3+2p^2-p)d+8p^5-2p^4.
\]
The leading coefficient is positive and
\[
 \operatorname{disc}_d(P_p)
 =-p^2(p-1)^2E(p)\leq0.
\]
Indeed, on writing $u=p-1/2\geq0$,
\[
 E(p)=92u^4+196u^3+135u^2+36u+2>0.
\]
Therefore $P_p(d)\geq0$.

For the remaining lower endpoint put
\[
 C_p(d)=G_p\bigl(d,d-(1-p)\bigr).
\]
It is needed only for $d\in[1-p/2,1]$.  Its leading coefficient as a cubic
in $d$ is $2(3p^3-p^2+1)>0$.  For $p<1$, exact expansion gives
\[
 \operatorname{disc}_d(C_p)
 =-8p^4(p-1)^4D(p)<0,
\]
where, for $u=p-1/2\geq0$,
\begin{align*}
 D(p)={}&736u^{10}+2624u^9+4448u^8+6756u^7+5532u^6
 +776u^5\\
 &+\frac{2331}{2}u^4+\frac{2161}{4}u^3
 +\frac{1185}{8}u^2+\frac{197}{8}u+\frac{51}{32}>0.
\end{align*}
Thus $C_p$ has exactly one real zero.  At the left endpoint,
\[
 C_p(1-p/2)
 =\frac{p^2(2-p)}4(29p^3-46p^2+24p-4)>0.
\]
The cubic in the last parentheses has discriminant $-304$, positive leading
coefficient, and value $1/8$ at $p=1/2$; its unique real zero is therefore
less than $1/2$.  It follows that the unique zero of $C_p$ lies below
$1-p/2$, and $C_p(d)>0$ throughout the required interval.  When $p=1$,
the same conclusion follows directly from
\[
 C_1(d)=6d(d-1)^2\geq0.
\]

It remains to check an interior stationary point.  Solving
$\partial G_p/\partial a=0$ gives
\[
 a_* =\frac p6(1+\rho),
 \qquad
 \rho=\sqrt{1+12d(3p-1)}.
\]
If $a_*$ belongs to the present region, then $a_*\geq p/2$ and hence
$\rho\geq2$.  A direct factorization gives
\[
 G_p(d,a_*)
 =\frac{p^2(\rho+1)(6p-1-\rho)^2}
 {864(3p-1)^2}\,R_p(\rho).
\]
Writing $r=\rho-2\geq0$ and $u=p-1/2\geq0$, the remaining factor is
\[
 R_p(\rho)
 =r^3+(9+12u)r^2+(24+66u+48u^2)r
  +14+50u+48u^2>0.
\]
All endpoint and stationary-point minima are nonnegative, completing the
proof.
\end{proof}

\begin{remark}[Why the density hypothesis is essential]
The coefficients in Lemma~\ref{lem:scalar}, the lower constraint
$a\geq d-(1-p)$, and the sign certificates all use $p\geq1/2$.  The lemma
is not an unconditional tree-domination statement.  This is exactly what
allows it to avoid the forest-domination obstruction recorded in
\texttt{FAILED\_PROOF\_ATTEMPTS.md}.
\end{remark}


%%%%%%%%%%%%%%%%%%%%%%%%%%%%%%%%%%%%%%%%%%%%%%%%%%%%%%%%%%%%%%%%%%%%%%%%%%%%%%%
\section{The positive graphon theorem}
%%%%%%%%%%%%%%%%%%%%%%%%%%%%%%%%%%%%%%%%%%%%%%%%%%%%%%%%%%%%%%%%%%%%%%%%%%%%%%%

\begin{theorem}[Triangle with a two-leaf broom]\label{thm:main}
For every graphon $W$ on an arbitrary probability space, of edge density
$p\in[1/2,1]$,
\[
 \boxed{t(B,W)\geq p^4(2p-1)=\Phi_B(p).}
\]
\end{theorem}

\begin{proof}
For $d>0$ define
\[
 h_p(d,a)=\frac{a^2}{d}(2a-p)_+,
\]
and set $h_p(0,0)=0$.  By \eqref{eq:density} and
Lemma~\ref{lem:pointwise},
\[
 t(B,W)\geq\int_\Omega h_p(d(x),A(x))\dd\mu(x).
\]
Lemma~\ref{lem:region} supplies the pointwise hypotheses of
Lemma~\ref{lem:scalar}.  Integrating that scalar inequality and using
\[
 \int d=p,
 \qquad
 \int(A-d^2)=0,
\]
gives
\begin{align*}
 t(B,W)
 &\geq2p^4(1-4p)+p^3(10p-3)p\\
 &=p^4(2p-1),
\end{align*}
which is \eqref{eq:target}.
\end{proof}

At $p=1/2$ the target is zero and the proof above remains valid.  At $p=1$,
one may alternatively note that $0\leq W\leq1$ and $\int W=1$ imply
$W=1$ almost everywhere, so both sides equal one.  No division by
$2p-1$ occurs.

For every integer $k\geq2$, the balanced complete $k$-partite graphon has
$p=1-1/k$ and realizes equality throughout the proof:
$d=p$, $A=p^2$, $C=p^3$, and
$\tau=p(2p-1)$ almost everywhere.  Thus the bound has the expected
chromatic extremizers.


%%%%%%%%%%%%%%%%%%%%%%%%%%%%%%%%%%%%%%%%%%%%%%%%%%%%%%%%%%%%%%%%%%%%%%%%%%%%%%%
\section{Atlas consequence, hidden-assumption audit, and limitations}
%%%%%%%%%%%%%%%%%%%%%%%%%%%%%%%%%%%%%%%%%%%%%%%%%%%%%%%%%%%%%%%%%%%%%%%%%%%%%%%

\begin{corollary}[Atlas 100]\label{cor:atlas100}
NetworkX Atlas graph $100$, with graph6 code \texttt{EsCW}, is positive.
It has order six, size six, chromatic number three, chromatic polynomial
\[
 \chi_{F_{100}}(x)=x(x-1)^4(x-2),
\]
and every graphon of edge density $p\in[1/2,1]$ satisfies
\[
 t(F_{100},W)\geq p^4(2p-1)=\Phi_{F_{100}}(p).
\]
\end{corollary}

\begin{proof}
The NetworkX edge list is
\[
 01,\ 02,\ 03,\ 34,\ 35,\ 45.
\]
Vertices $3,4,5$ form a triangle, rooted at $3$.  Vertex $0$ is joined to
that root, and $1,2$ are precisely the two leaves adjacent to $0$.  Hence
the graph is exactly $B$, and Theorem~\ref{thm:main} applies.
\end{proof}

All graphon integrands used above are bounded, nonnegative, and measurable,
so Tonelli--Fubini applies on an arbitrary probability space.  Every
pointwise statement is understood almost everywhere.  The proof uses
ordinary homomorphism densities, not injective copies.  It assumes neither
a finite-step representation, regularity, constant degree, nor a positive
minimum degree.  The set $\{d=0\}$ is handled without division.

For formal verification, the reusable new item is
Lemma~\ref{lem:scalar}.  A convenient implementation can split on
$d=0$, $a\leq p/2$, and $a\geq p/2$, exactly as above.  The graphon layer
then needs only Cauchy--Schwarz under the measure $W(x,y)\dd\mu(y)$, the
already used rooted-triangle inequality, the two identities of
Lemma~\ref{lem:region}, and integration of the scalar supporting plane.
The graph-specific module awaiting inspection is
\texttt{lean/Taeyoung/Examples/Graph100.lean}.

The theorem covers exactly a triangle with one adjacent star centre and two
leaves at that centre.  It does not prove arbitrary pendant-tree closure,
a star with another number of leaves, several such brooms, or brooms rooted
at different triangle vertices.  The supporting plane is special to the
two-leaf exponent and the interval $p\geq1/2$; none of these extensions may
be inferred without a new scalar certificate.

Every polynomial identity, discriminant, positive coefficient shift, the
Atlas isomorphism, and the chromatic polynomial used in this document are
independently checked over exact rational arithmetic by
\begin{verbatim}
python codes/verify_triangle_two_leaf_broom.py
\end{verbatim}
This audit is only a check of the displayed algebra; the arbitrary-graphon
argument is the proof given above.

%%%%%%%%%%%%%%%%%%%%%%%%%%%%%%%%%%%%%%%%%%%%%%%%%%%%%%%%%%%%%%%%%%%%%%%%%%%%%%%
\section{What the Lean formalization actually proves}
%%%%%%%%%%%%%%%%%%%%%%%%%%%%%%%%%%%%%%%%%%%%%%%%%%%%%%%%%%%%%%%%%%%%%%%%%%%%%%%

The machine-checked development is
\texttt{lean/Taeyoung/Methods/Broom.lean}, ending at
\texttt{satisfiesLowerBound\_100}.  Theorem~\ref{thm:main} is proved in full, on
all of $p\in[1/2,1]$.  The departure is inside Lemma~\ref{lem:scalar}, in its
active case $a\geq p/2$.

\paragraph{The departure: no cubic discriminant.}
This note treats $G_p(d,a)=2a^3-pa^2-dL_p(d,a)$ as a cubic in $a$ for fixed $d$
and runs a case analysis on its root structure.  The formalization does it in
one chain.  In the coordinates
\[
 \alpha=a-p^2,\qquad \delta=d-p
\]
--- the point where the extremal graphon sits --- the cleared inequality is the
\texttt{ring} identity
\begin{equation}
 a^2(2a-p)-d\,L_p(d,a)
 =p(6p-1)\alpha^2-2p^2(3p-1)\alpha\delta+p^3(8p-3)\delta^2
 +2\alpha^3+2p^2(3p-1)\delta^3 .
 \label{eq:lean-broom}
\end{equation}
Both cubic terms are absorbed using only the two facts that already hold in the
active case --- $a\geq p/2$ by assumption and $d\geq a\geq p/2$:
\[
 2\alpha^3\geq(p-2p^2)\alpha^2,
 \qquad
 2p^2(3p-1)\delta^3\geq-p^3(3p-1)\delta^2 ,
\]
each being $2\alpha^2(a-\frac p2)\geq0$ and
$2p^2(3p-1)\delta^2(d-\frac p2)\geq0$ rearranged.  What remains,
\[
 4p^2\alpha^2-2p^2(3p-1)\alpha\delta+p^3(5p-2)\delta^2 ,
\]
is positive definite on $p\geq1/2$, certified by
\begin{equation}
 16p^2\bigl(4p^2\alpha^2-2p^2(3p-1)\alpha\delta+p^3(5p-2)\delta^2\bigr)
 =\bigl(8p^2\alpha-2p^2(3p-1)\delta\bigr)^2+4p^4(11p^2-2p-1)\delta^2,
 \label{eq:lean-broom-pd}
\end{equation}
with $11p^2-2p-1\geq3/4$ there.  So the active case is four \texttt{ring}
identities and three sign checks, and no cubic is analysed.

The inactive case $a\leq p/2$ follows this note exactly, including its two
subcases $d\leq p/2$ (push $a$ up to $d$, then $Q_1\geq Q_1(p/2)>0$) and
$d\geq p/2$ (push $a$ up to $p/2$, then $Q_2\geq0$ by its nonpositive
discriminant $-p(2p-1)^2(23p-8)$).

\paragraph{Everything else transcribes directly.}
Lemma~\ref{lem:region} is \texttt{Broom.le\_pathOp} together with
\texttt{K4Tail.pathOp\_le\_degree} and \texttt{BookTail.integral\_pathOp};
Lemma~\ref{lem:pointwise} is \texttt{Broom.sq\_pathOp\_le} (the project's
weighted Cauchy--Schwarz at weight $W(x,\cdot)$) and the pre-existing
\texttt{Link.rooted\-Triangle\_ge}.  The division by $d$ is never performed:
\texttt{Broom.mul\_plane\_le} states the pointwise bound multiplied by $d$, so
the zero-degree set is an ordinary case rather than a convention.  The
identity \eqref{eq:density} is \texttt{Broom.homDensity\_broom}, through the
generic peeling ladder of \texttt{Methods/Peeling.lean}.


\end{document}
